\section{Conclusion}
\label{sec:conclusion}

We presented a controlled benchmark of hidden-state similarity in plain four-layer MLPs trained on \mnist, \fashionmnist, and \cifar. By combining linear \cka, \svcca, \pwcca, cosine similarity, and per-layer linear probes, we separated within-model layer locality from cross-seed recurrence and from task-information growth with depth.

Our main finding is simple: MLP hidden spaces across depth are not totally different. Adjacent layers are consistently more similar than distant ones, while later layers usually improve decodable label information. Cross-seed comparisons are weaker, which shows that representational continuity within a model is easier to establish than stable layer identity across random initializations.

The next steps are straightforward. Functional stitching can test whether these similarities support interchangeable computation. Additional sweeps over width, depth, and pre- versus post-activation representations can identify when progressive similarity stops holding. Extending the same protocol to residual networks and transformers would clarify how architecture changes the geometry of layer-wise progression.
